\section{Requirements Analysis, Specification and Validation}
\label{sec:pa-analysis}

\subsection{Requirements Analysis}
\label{sec:pa-analyst}
\todo{The analyst extracts raw requirements (each \texttt{stated\_by} a persona and \texttt{evidenced\_by} an interview message), consolidates them into analyzed requirements with a priority and a quality category, and records latent needs, assumptions, ambiguity issues and conflicts. When open issues need a stakeholder's answer, the result is ``not OK'' and a rework directive is written (Section~\ref{sec:pa-deliberation}).}

\subsection{Requirement Specification}
\label{sec:pa-specification}
\todo{Structure of the specification (sections, functional and non-functional requirements); each specification requirement \texttt{derived\_from} analyzed requirements; quality per ISO/IEC/IEEE 29148~\cite{iso29148}.}

\subsection{Product Backlog}
\label{sec:pa-backlog}
\todo{User stories with acceptance criteria and priority, each \texttt{derived\_from} a specification requirement; INVEST~\cite{cohn2004user}; dependencies.}

\subsection{Validation}
\label{sec:pa-validation}
\todo{Validation categories; a validation issue \texttt{flags} the items it is about; on fail the run returns to the requirement engineer.}
